\documentclass[11pt]{article}

\usepackage[letterpaper,margin=1in]{geometry}
\usepackage[T1]{fontenc}
\usepackage[utf8]{inputenc}
\usepackage{lmodern}
\usepackage{microtype}
\usepackage{amsmath,amssymb,mathtools,bm}
\usepackage{xcolor}
\usepackage{enumitem}
\usepackage{hyperref}
\usepackage[most]{tcolorbox}
\usepackage{fancyhdr}
\usepackage{parskip}

\definecolor{qftblue}{HTML}{245A8D}
\definecolor{qftgray}{HTML}{F4F6F8}
\hypersetup{
  colorlinks=true,
  linkcolor=qftblue,
  urlcolor=qftblue,
  pdftitle={PHYS 662 Homework Set 3},
  pdfauthor={Nima Lashkari}
}

\pagestyle{fancy}
\fancyhf{}
\lhead{PHYS 662: Quantum Field Theory I}
\rhead{Homework Set 3}
\cfoot{\thepage}
\setlength{\headheight}{14pt}

\setlist[enumerate]{leftmargin=2.2em,itemsep=0.55em,topsep=0.5em}
\newlist{parts}{enumerate}{1}
\setlist[parts]{label=\textbf{(\alph*)},leftmargin=2.3em,itemsep=0.65em}

\newcommand{\dd}{\mathrm{d}}
\newcommand{\ii}{\mathrm{i}}
\newcommand{\R}{\mathbb{R}}
\newcommand{\C}{\mathbb{C}}
\newcommand{\Z}{\mathbb{Z}}
\newcommand{\Lag}{\mathcal L}
\newcommand{\tr}{\operatorname{tr}}
\newcommand{\Ad}{\operatorname{Ad}}

\newtcolorbox{conventionbox}[1][]{
  enhanced,breakable,
  colback=qftgray,colframe=qftblue,
  boxrule=0.8pt,arc=1.5mm,
  title={#1},fonttitle=\bfseries,coltitle=white
}
\newcommand{\problem}[2]{%
  \clearpage
  \section*{\color{qftblue}Problem #1: #2}
  \addcontentsline{toc}{section}{Problem #1: #2}
}
\newcommand{\lectureconnection}[1]{%
  {\small\textit{Lecture connection: #1.}\par}\medskip
}

\begin{document}

\begin{center}
  {\scshape\Large PHYS 662 --- Quantum Field Theory I\par}
  \vspace{0.25em}
  {\Huge\bfseries\color{qftblue} Homework Set 3\par}
  \vspace{0.25em}
  {\large Lectures 7--10 \quad Fall 2026\par}
\end{center}

\vspace{0.5em}
\noindent
\textbf{Name:} \rule{0.44\textwidth}{0.4pt}
\hfill
\textbf{Due:} \rule{0.24\textwidth}{0.4pt}

\begin{conventionbox}[Instructions and conventions]
Show the reasoning leading to each result. Use \(\hbar=c=1\),
\(\eta_{\mu\nu}=\operatorname{diag}(1,-1,-1,-1)\), and
\(\epsilon_{123}=+1\). Work in \(3+1\) dimensions unless stated otherwise.
Representation spaces are complex unless a real space is specified.

For Hermitian generators, use
\[
 U(\theta)=e^{-\ii\theta^a t_a},\qquad
 [t_a,t_b]=\ii f_{ab}{}^c t_c.
\]
Repeated indices are summed. Distinguish field-component indices from
generator labels. Classical transformations are generated by
\(\delta F=\epsilon^a\{F,Q_a\}_{\rm PB}\); the corresponding quantum
action is \(\delta\widehat F=\ii\epsilon^a[\widehat Q_a,\widehat F]\).

Keep boundary terms until their role is established. For fields on
\(\R^3\), assume sufficient spatial falloff for the stated charges and
moments to exist and for the relevant surface terms to vanish.
Treat complex fields and their conjugates as independent variables when
varying the action. Distinguish identities valid off shell from
conservation laws that use the equations of motion.
\end{conventionbox}

The six problems are weighted equally.

\tableofcontents

\problem{1}{Algebra of conserved charges}
\lectureconnection{7 and 10}

Let \(\Phi\) be a column of \(N\) complex scalar fields with
\begin{equation}
 \Lag=\partial_\mu\Phi^\dagger\partial^\mu\Phi
       -V(\Phi,\Phi^\dagger).
 \label{eq:multiplet}
\end{equation}
The differentiable real potential is invariant under a compact internal
group \(G\subseteq U(N)\), with constant transformations
\(\Phi'=e^{-\ii\theta^a t_a}\Phi\).
For example, \(V=V(\Phi^\dagger\Phi)\) has a \(U(N)\) symmetry.

\begin{parts}
 \item Derive the field equations and the Noether currents \(j_a^\mu\).
 Verify \(\partial_\mu j_a^\mu=0\) using the field equations and the
 infinitesimal invariance of \(V\). Define the charges \(Q_a\) and state
 the boundary condition needed for their conservation.

 \item Write the canonical momenta
 \(\Pi_A=\dot\Phi_A^\dagger\) and
 \(\overline\Pi^{\,A}=\dot\Phi^A\), and express \(Q_a\) in these
 variables. Quantize with the equal-time commutators
 \begin{align}
 [\widehat\Phi^A(t,\bm x),\widehat\Pi_B(t,\bm y)]
 &=\ii\delta^A{}_B\delta^{(3)}(\bm x-\bm y),\notag\\
 [\widehat\Phi_A^\dagger(t,\bm x),
       \widehat{\overline\Pi}^{\,B}(t,\bm y)]
 &=\ii\delta_A{}^B\delta^{(3)}(\bm x-\bm y),
 \label{eq:fieldCCR}
 \end{align}
 with all other elementary commutators zero. Compute
 \([\widehat Q_a,\widehat\Phi]\) and
 \([\widehat Q_a,\widehat\Phi^\dagger]\), and check that they generate
 the stated transformations, including their signs.

 \item Let \(\widehat\rho_a=\widehat j_a^0\). Calculate explicitly
 \begin{equation}
 [\widehat\rho_a(t,\bm x),\widehat\rho_b(t,\bm y)],\qquad
 [\widehat Q_a,\widehat\rho_b(t,\bm x)],\qquad
 [\widehat Q_a,\widehat Q_b].
 \label{eq:requestedAlgebra}
 \end{equation}
 Express the results in terms of \(f_{ab}{}^c\), and track the spatial
 delta function before integrating. Use a spatial regulator preserving
 \(G\) and symmetric ordering of the canonical bilinears; compute this
 canonical algebra before taking its continuum limit. Compare with the
 classical Poisson brackets. Explain why conservation alone would not
 establish the charge algebra or exclude a central term.

 \item Specialize to \(G=U(N)\) and \(V=V(\Phi^\dagger\Phi)\), with
 \(\tr(t_at_b)=\tfrac12\delta_{ab}\) and
 \(t_0=I_N/\sqrt{2N}\). Count the \(U(N)\) currents and the currents in
 its \(SU(N)\) subgroup. Does \(\widehat Q_0\) commute with every
 \(\widehat Q_a\)? Does it commute with \(\widehat\Phi\)? Explain how
 the answers are compatible.
\end{parts}

\problem{2}{The Galilean group and its mass extension}
\lectureconnection{7, 8, and 10}

Consider \(n\) spinless nonrelativistic particles with masses \(m_r>0\),
canonical brackets \(\{x_r^i,p_s^j\}_{\rm PB}=\delta_{rs}\delta^{ij}\),
and Lagrangian
\begin{equation}
 L=\sum_{r=1}^{n}\frac{m_r}{2}\dot{\bm x}_r^{\,2}
   -\sum_{r<s}W_{rs}(|\bm x_r-\bm x_s|).
 \label{eq:galileanL}
\end{equation}
The differentiable pair potentials have no explicit time dependence.
Include spatial translations, time translations, rotations, and boosts.
Use the boost convention \(\delta\bm x_r=-\bm v t\), with time unchanged.

\begin{parts}
 \item Find the Hamiltonian \(H\), total momentum \(\bm P\), and orbital
 angular momentum \(\bm J\). Under a boost, show that the first-order
 change of \(L\) is a total time derivative. Retain this boundary term
 in Noether's theorem and derive the boost charge \(\bm K\).
 Check both its canonical action on \(\bm x_r,\bm p_r\) and its
 conservation using
 \begin{equation}
 \frac{\dd K_i}{\dd t}
 =\frac{\partial K_i}{\partial t}+\{K_i,H\}_{\rm PB}.
 \label{eq:explicitTime}
 \end{equation}

 \item Compute the Poisson brackets among \(H,P_i,J_i,K_i\), displaying
 all nonzero brackets and identifying the vanishing ones. In particular,
 evaluate \(\{K_i,P_j\}_{\rm PB}\), \(\{K_i,H\}_{\rm PB}\), and
 \(\{K_i,K_j\}_{\rm PB}\). Identify the role of
 \(M=\sum_r m_r\), and explain in what sense it is central.

 \item Quantize the generators and write the centrally extended
 commutator algebra. In the abstract algebra, regard \(\widehat M\) as
 an additional generator; in this fixed-mass system it is represented by
 \(M I\). Check the Jacobi identities involving boosts, translations,
 and time translations, and the compatibility of the central term with
 rotations. Show that constant shifts of the generators cannot remove
 the mass term. Did this extension first appear upon quantization?

 \item Define
 \begin{equation}
 B(\bm v)=e^{-\ii\bm v\cdot\widehat{\bm K}},\qquad
 T(\bm a)=e^{-\ii\bm a\cdot\widehat{\bm P}}.
 \end{equation}
 Use the Baker--Campbell--Hausdorff formula to evaluate
 \(B(\bm v)T(\bm a)B(\bm v)^{-1}T(\bm a)^{-1}\).
 Compare with the composition of the corresponding spacetime maps.
 Explain why the phase gives a projective representation of the
 unextended group and an ordinary representation of its central extension.
\end{parts}

\problem{3}{The adjoint representation and the Jacobi identity}
\lectureconnection{8 and 9}

Let \(G\) be a connected matrix Lie group with an \(r\)-dimensional
Lie algebra. Use a basis \(T_a\) in the convention
\begin{equation}
 [T_a,T_b]=\ii f_{ab}{}^cT_c,
 \qquad a,b,c=1,\ldots,r.
 \label{eq:abstractAlgebra}
\end{equation}
Work on the complexified span of these generators when writing
representation matrices. Define
\begin{equation}
 gT_bg^{-1}=T_cD(g)^c{}_b,\qquad
 (t_a^{\rm ad})^c{}_b=\ii f_{ab}{}^c.
 \label{eq:adjointMatrices}
\end{equation}

\begin{parts}
 \item Prove \(D(g_1g_2)=D(g_1)D(g_2)\). For
 \(g=e^{-\ii\theta^aT_a}\), expand conjugation to first order and
 verify that \(D(g)=I-\ii\theta^at_a^{\rm ad}+O(\theta^2)\).
 State which index of \((t_a^{\rm ad})^c{}_b\) is the row index.

 \item Compute the matrix commutator in components and show that
 \begin{equation}
 [t_a^{\rm ad},t_b^{\rm ad}]
 =\ii f_{ab}{}^c t_c^{\rm ad}.
 \label{eq:adjointAlgebra}
 \end{equation}
 Write the resulting identity among structure constants. Relate it to
 \begin{equation}
 [T_a,[T_b,T_c]]+[T_b,[T_c,T_a]]+[T_c,[T_a,T_b]]=0,
 \end{equation}
 and explain why this identity follows from associativity of matrix
 multiplication.

 \item Determine the dimension of the adjoint representation.
 Show that its invariant subspaces correspond to ideals of the
 complexified Lie algebra. Identify the kernel of the infinitesimal
 adjoint action in terms of the center of the Lie algebra. Distinguish
 this kernel from the discrete center that a connected group may have.

 \item For \(\mathfrak{su}(2)\), use \(f_{ab}{}^c=\epsilon_{abc}\)
 to construct the three \(3\times3\) matrices \(t_a^{\rm ad}\).
 Find the eigenvalues of \(t_3^{\rm ad}\) and identify the spin.
 Then give the dimensions of the adjoint representations of
 \(\mathfrak{su}(N)\) and \(\mathfrak u(N)\), for \(N\ge2\).
 Explain the decomposition of the latter into an invariant singlet and
 a traceless part. Does the central \(U(1)\) act nontrivially by
 conjugation?
\end{parts}

\problem{4}{The topology of rotations and half-integer spin}
\lectureconnection{9 and 10}

Establish \(\pi_1(SO(3))\cong\Z_2\) and explain why \(SO(3)\) is
not simply connected whereas \(SU(2)\) is. Use the Pauli-matrix identity
\begin{equation}
 \sigma_i\sigma_j=\delta_{ij}I_2+\ii\epsilon_{ijk}\sigma_k.
 \label{eq:pauli}
\end{equation}
You may use that \(S^3\) is simply connected and that paths and homotopies
lift uniquely through a covering map once their initial lift is fixed.

\begin{parts}
 \item Show that every \(U\in SU(2)\) can be written uniquely as
 \begin{equation}
 U=a_0I_2-\ii\bm a\cdot\bm\sigma,
 \qquad a_0,\bm a\in\R,\qquad a_0^2+|\bm a|^2=1.
 \label{eq:su2Sphere}
 \end{equation}
 Explain how this identifies the group manifold with \(S^3\).

 \item Define a real linear map \(R_U\) by
 \begin{equation}
 U(\bm x\cdot\bm\sigma)U^\dagger
 =(R_U\bm x)\cdot\bm\sigma.
 \label{eq:cover}
 \end{equation}
 Show that \(R_U\in SO(3)\), that \(R_{UV}=R_UR_V\), and that
 \(U=e^{-\ii\theta\bm n\cdot\bm\sigma/2}\) produces a rotation by
 \(\theta\) about \(\bm n\). Determine the kernel, establish
 surjectivity, and explain why this map is a two-sheeted covering.

 \item A loop in \(SO(3)\), based at the identity, lifts to a path in
 \(SU(2)\) starting at \(I_2\). Determine its possible endpoints and
 use them to classify loops up to based homotopy. Explain why there are
 exactly two classes and how concatenation gives the group \(\Z_2\).
 Apply this argument to a \(2\pi\) rotation and to a \(4\pi\) rotation
 about a fixed axis. A full turn returning a classical vector to itself
 is not, by itself, a proof that the corresponding loop is contractible.

 \item In the spin-\(j\) representation, use the eigenvalues
 \(m=-j,-j+1,\ldots,j\) of \(J_3\) to evaluate
 \(e^{-2\pi\ii J_3}\). Which representations descend to ordinary
 representations of \(SO(3)\), and which give projective
 representations on quantum rays? Explain why half-integer spin does
 not require a nonzero central term in
 \([J_i,J_j]=\ii\epsilon_{ijk}J_k\).
\end{parts}

\problem{5}{Spacetime symmetries and the center of energy}
\lectureconnection{7}

Consider a real scalar field on Minkowski spacetime,
\begin{equation}
 \Lag=\frac12\partial_\mu\phi\,\partial^\mu\phi-V(\phi),
 \qquad V(\phi)\ge0,\qquad V(0)=0,
 \label{eq:realScalar}
\end{equation}
where \(V\) is differentiable. Assume localized configurations with
sufficient falloff for the energy and its first spatial moment to be
finite. Use the active convention \(x'=x+\xi(x)\).

\begin{parts}
 \item For a transported scalar, distinguish
 \(\Delta\phi=\phi'(x')-\phi(x)\) from
 \(\delta\phi=\phi'(x)-\phi(x)\). For translations
 \(\xi^\mu=a^\mu\), derive the Noether current, including its sign,
 and identify the canonical stress tensor \(T^\mu{}_{\nu}\).
 Verify its conservation directly from the field equation. Define
 the conventional physical momentum by
 \(P^\nu=\int\dd^3x\,T^{0\nu}\), and relate it to the direct Noether
 charge for the parameter \(a^\mu\).

 \item Compute the energy density \(\mathcal H=T^{00}\), momentum
 density \(T^{0i}\), and energy flux \(T^{i0}\). Check the energy density
 by a Legendre transform and write its local continuity equation.
 Explain which statements require the field equation and which follow
 directly from the form of the Lagrangian.

 \item Derive the Lorentz current for
 \(\xi^\mu=\omega^\mu{}_{\nu}x^\nu\), with
 \(\omega_{\mu\nu}=-\omega_{\nu\mu}\), and verify its conservation.
 Express the spatial rotation charges in terms of the momentum density.
 Show that the boost charges can be written as
 \begin{equation}
 B^i=tP^i-\int\dd^3x\,x^i\mathcal H.
 \label{eq:boostCharge}
 \end{equation}
 For \(H=\int\mathcal H\,\dd^3x>0\), define the center of energy
 \(X_E^i=H^{-1}\int x^i\mathcal H\,\dd^3x\). Derive its velocity and
 explain why the explicit time dependence of \(B^i\) is compatible
 with conservation.

 \item Now add a prescribed source,
 \(\Lag_h=\Lag+h(x)\phi\). In testing an active symmetry, hold the
 background function \(h(x)\) fixed. Derive the field equation and the
 on-shell divergence of the canonical tensor \(T_h^\mu{}_{\nu}\),
 including its sign. When is energy conserved? Which component of
 momentum remains conserved if \(h\) is independent of one spatial
 coordinate? Explain why a coordinate relabeling alone cannot restore
 a broken translation symmetry of this fixed-background system.
\end{parts}

\problem{6}{Two spins, invariant subspaces, and exchange}
\lectureconnection{8 and 9}

Let \(V_{1/2}\cong\C^2\) carry the spin-one-half representation,
with \(S_i=\sigma_i/2\). On \(V_{1/2}\otimes V_{1/2}\), define
\begin{equation}
 S_{1i}=S_i\otimes I_2,\qquad S_{2i}=I_2\otimes S_i,\qquad
 J_i=S_{1i}+S_{2i}.
 \label{eq:twoSpins}
\end{equation}
The two factors are labeled spin spaces; impose no identical-particle
constraint. Use the product basis built from
\(S_3|\uparrow\rangle=\tfrac12|\uparrow\rangle\) and
\(S_3|\downarrow\rangle=-\tfrac12|\downarrow\rangle\).

\begin{parts}
 \item Differentiate \(U\otimes U\) to obtain \(J_i\), and verify
 \([J_i,J_j]=\ii\epsilon_{ijk}J_k\).
 Explain the distinction between the dimension of \(SU(2)\), the
 dimension of \(V_{1/2}\), and the dimension of the tensor-product space.

 \item Starting with the highest-weight state and using
 \(J_-=J_1-\ii J_2\), construct a normalized triplet. Find the normalized
 state orthogonal to the triplet and show that every \(J_i\) annihilates
 it. Give the \(J^2\) and \(J_3\) eigenvalues of all four states, and
 establish
 \begin{equation}
 V_{1/2}\otimes V_{1/2}\cong V_0\oplus V_1.
 \label{eq:spinDecomposition}
 \end{equation}

 \item Derive the singlet and triplet projectors \(P_0,P_1\) as
 polynomials in \(\bm S_1\cdot\bm S_2\). Verify that they are
 orthogonal projectors and sum to the identity. Use Schur's lemma to
 find the most general Hermitian Hamiltonian commuting with every
 \(U\otimes U\). Express it both in terms of \(P_0,P_1\) and in terms
 of \(I_4,\bm S_1\cdot\bm S_2\), and relate the coefficients to its
 singlet and triplet energies.

 \item Define the exchange operator
 \(\mathsf X(v\otimes w)=w\otimes v\). Determine its eigenvalues
 on the two invariant subspaces and express \(\mathsf X\) in terms
 of \(\bm S_1\cdot\bm S_2\). Does exchange symmetry alone force a
 Hamiltonian to be rotationally invariant? Justify your answer with
 an argument or a counterexample.

 \item Evaluate the action of the central element \(-I_2\in SU(2)\)
 on the tensor product. Does this representation descend to an ordinary
 representation of \(SO(3)\)? Compare with the direct sum
 \(V_0\oplus V_{1/2}\): does that space admit even projective descent
 on all its rays? State the condition on the action of the covering
 kernel that distinguishes the two cases.
\end{parts}

\end{document}
